\documentclass[11pt]{article}
\usepackage[margin=1in]{geometry}
\usepackage{amsmath,amssymb,amsthm,hyperref}
\hypersetup{colorlinks=true,urlcolor=blue,linkcolor=blue,citecolor=blue}
\newtheorem{theorem}{Theorem}
\newtheorem{lemma}{Lemma}
\newtheorem{corollary}{Corollary}
\newcommand{\Av}{\operatorname{Av}}
\title{Marked primitive components and the expansion conjecture\\
\large MathDB 360957}
\author{Principia Math in Progress}
\date{September 18, 2026}
\begin{document}
\maketitle

\section*{Task 1. Full solution}
For $a,k\ge1$, let $\mathcal S^{a\prec n}_{n,k}(1324)$ consist of
the permutations of $[n]$ avoiding $1324$ in which $a$ is $k$ positions
to the left of $n$, and all letters $1,\ldots,a-1$ are to the right
of $n$. Write
\[
 T_{a,k}(x)=\sum_n |\mathcal S^{a\prec n}_{n,k}(1324)|x^n,
 \qquad T_{a,0}(x)=|\Av_{a-1}(1324)|x^a.
\]
The auxiliary class at $k=0$ consists of the words $a\lambda$, where
$\lambda$ is a $1324$-avoiding permutation of $[a-1]$. The empty
permutation is allowed, so $T_{1,0}=x$.

A \emph{primitive} is a $1324$-avoiding permutation whose minimum
is immediately followed by its maximum. If $\mathcal P_m$ is this
class in size $m\ge2$, the series used in the source is
\[
 f(x)=\sum_{m\ge2}|\mathcal P_m|x^{m-1}
     =x+2x^2+6x^3+22x^4+91x^5+\cdots.
\]
We use the primitive factorization of Gil--Lopez--Weiner
\cite[Propositions 2.3--2.4 and Corollary 2.5]{GLW}.
Their paper proves the cases $a=1,2$ and also gives the marking
argument for $a=3$. The following proof extends that method to
every $a$ by proving that unmarked components can always be inserted.

\begin{theorem}\label{main}
For each $a\ge1$ there are series $H_{a,r}(x)$ with nonnegative
integer coefficients, zero for $r\ge a$, such that for every $k\ge0$,
\begin{equation}\label{expansion}
 T_{a,k}=\sum_{r=0}^{\min(k,a-1)}\binom{k}{r}f^{k-r}H_{a,r}.
\end{equation}
Consequently, for every $k\ge a$ all three formulas of
\cite[Conjecture 1]{GLW} hold:
\begin{align}
 &\sum_{j=0}^k(-1)^j\binom{k}{j}f^jT_{a,k-j}=0,\label{first}\\
 &T_{a,k}=\sum_{r=0}^{a-1}\binom{k}{r}f^{k-r}
       \sum_{i=0}^r(-1)^i\binom{r}{i}f^iT_{a,r-i},\label{second}\\
 &T_{a,k}=\sum_{j=0}^{a-1}(-1)^{a-j-1}\binom{k}{j}
       \binom{k-j-1}{a-j-1}f^{k-j}T_{a,j}.\label{third}
\end{align}
All identities are identities of formal power series in $x$.
\end{theorem}

\paragraph{The primitive components.}
Call letters below $a$ \emph{small}, and letters at least $a$
\emph{large}. Delete the small letters from an object counted by
$T_{a,k}$ and standardize. The cited unique factorization gives
$k$ primitives. Before standardizing, this large-letter word has
the form
\begin{equation}\label{skeleton}
 A_k\cdots A_1\;a\,b_1\cdots b_k\;B_k\cdots B_1,
 \qquad a=b_0<b_1<\cdots<b_k=n.
\end{equation}
For each $i$, the words $A_i,B_i$ partition the interval of integer
letters $(b_{i-1},b_i)$. The word
$A_i\,b_{i-1}\,b_i\,B_i$ standardizes to the $i$th primitive.
Conversely, any ordered list of primitives gives a unique avoiding
word of the form \eqref{skeleton}. In particular, every large
letter in $B_i$ exceeds every large letter in $B_j$ when $i>j$.
The suffix $B_k\cdots B_1$ avoids $213$: otherwise prepending the
earlier letter $a$ would give $1324$.

\paragraph{A canonical assignment of the small letters.}
In the suffix to the right of $n$, attach each maximal run of small
letters that is followed by a large letter to that following letter.
The final run, if any, is kept separately as a tail $L$.
Let $D_i$ be $B_i$ with its attached runs inserted. Thus the full
permutation has the unique form
\begin{equation}\label{decorated}
 A_k\cdots A_1\;a\,b_1\cdots b_k\;D_k\cdots D_1\;L.
\end{equation}
Every nonempty $D_i$ ends in a large letter, and $D_i$ is empty
if $B_i$ is empty. Mark component $i$ precisely when $D_i$ contains
a small letter. The marked components use disjoint nonempty sets
of the $a-1$ small letters; hence their number $r$ is at most $a-1$.
An unmarked component has $D_i=B_i$.

\begin{lemma}[Deletion and arbitrary reinsertion]\label{insertion}
Delete all unmarked components in \eqref{decorated}, meaning that
for each such $i$ we delete $A_i,B_i$, and the upper endpoint $b_i$.
Standardize the result, leaving the small letters $1,\ldots,a-1$
and the base letter $a$ fixed. The result is an avoiding permutation
with exactly $r$ components, all marked, or an auxiliary $k=0$
object if $r=0$.

Conversely, into any such fully marked object one may insert
arbitrary unmarked primitives at arbitrary places in its ordered
component list. This always produces a $1324$-avoiding permutation.
\end{lemma}

\begin{proof}
Deletion produces a subsequence followed by an order-preserving
relabeling, so it preserves avoidance. Each retained component
keeps its internal order and decorations. Its width
$b_i-b_{i-1}$ is unchanged; the successive intervals are simply
made consecutive again. The unique primitive factorization of the
remaining large-letter word is therefore exactly the retained list.
Each small run still has its assigned following large letter,
and the tail stays at the end. This proves the deletion assertion,
including $r=0$, when only $aL$ remains.

For the converse, form \eqref{decorated} with the extra primitives,
using consecutive intervals for their large letters and preserving
all small letters and their assignments. Its large-letter
subsequence avoids $1324$ by the cited primitive-product theorem.
Its small-letter subsequence is unchanged. Suppose nevertheless
that a forbidden pattern appears. Classify it by its number of
small letters.

With zero small letters it contradicts avoidance of the large-letter
subsequence. With four small letters it was already in the core.
With at least one small letter the first letter of the pattern
must be small: the first letter is the minimum of a $1324$
pattern. Thus the entire pattern lies to the right of the new
maximum, in $D_k\cdots D_1L$.

With exactly one small letter, the three large letters would form
$213$ in $B_k\cdots B_1$, which is impossible.

With exactly two small letters, the pattern has the form
$s_1,u,s_2,v$, where $s_1<s_2<a\le u<v$. An increasing pair
$u,v$ cannot cross two suffix blocks, since those blocks occur
in decreasing order of their value intervals. Hence $u,v$ lie
in the same $B_i$. This cannot be an inserted unmarked block,
because $s_2$ lies between $u$ and $v$. It is therefore a retained
marked block, and the same four letters form a forbidden pattern
in the core.

With exactly three small letters, those letters form $132$ and
are followed by a large letter $v$. If $v$ is retained, this
pattern is already in the core. If $v$ belongs to an inserted
block, consider the last of the three selected small letters,
in positional order. It cannot lie in the final tail. It lies
in a retained decorated block $D_i$ before the inserted block.
Its attached run is followed by a retained large letter $w$ in
that same $D_i$, hence after all three selected small letters
and before $v$. Replacing $v$ by $w$ gives $1324$ in the core,
again a contradiction. These cases exhaust all possibilities.
\end{proof}

\begin{proof}[Proof of Theorem \ref{main}]
Let $H_{a,r}$ count, by their full size, the avoiding objects with
$r$ components all marked, retaining every small letter and the
tail. For $r=0$ this is the auxiliary class, so $H_{a,0}=T_{a,0}$.
The marking bound gives $H_{a,r}=0$ for $r\ge a$.

Lemma \ref{insertion} is a bijection: an object with $k$ components
determines its core, the $r$ marked positions in the component
list, and its ordered $k-r$ unmarked primitives. Conversely these
data recover exactly one object. No labels need to be chosen:
consecutive large-letter intervals determine them, while the
small labels belong to the core. A primitive of size $m$ adds
$m-1$ letters, since its lower endpoint is shared with the
preceding part of the increasing central sequence. Thus every
unmarked primitive contributes $f$, and the choices of marked
positions contribute $\binom{k}{r}$. This proves \eqref{expansion}.

Ordinary weighted binomial inversion gives
\begin{equation}\label{inverse}
 H_{a,r}=\sum_{i=0}^r(-1)^i\binom{r}{i}f^iT_{a,r-i}.
\end{equation}
To check this, substitute \eqref{expansion} on the right.
The coefficient of $f^{r-s}H_{a,s}$ is
\[
 \binom{r}{s}\sum_{i=0}^{r-s}(-1)^i\binom{r-s}{i}.
\]
This is one if $s=r$ and zero otherwise. No division by $f$ is used.
Equation \eqref{first} follows because $H_{a,k}=0$ for $k\ge a$;
substitution into \eqref{expansion} gives \eqref{second}.

Finally the coefficient of $f^{k-j}T_{a,j}$ in \eqref{second} is
\begin{align*}
 \sum_{r=j}^{a-1}(-1)^{r-j}\binom{k}{r}\binom{r}{j}
 &=\binom{k}{j}\sum_{u=0}^{a-j-1}(-1)^u\binom{k-j}{u}\\
 &=(-1)^{a-j-1}\binom{k}{j}\binom{k-j-1}{a-j-1}.
\end{align*}
The last equality is the alternating partial-binomial identity
(an immediate telescoping consequence of Pascal's rule), valid
here because $k\ge a$. This proves \eqref{third}.
\end{proof}

\begin{corollary}
The binomial transforms \eqref{inverse} have nonnegative integer
coefficients. Moreover the positional generating function satisfies
\[
 g_a(x,t)=\sum_{r=0}^{a-1}\frac{H_{a,r}(x)t^r}
                                  {(1-tf(x))^{r+1}}-T_{a,0}(x).
\]
\end{corollary}
\begin{proof}
The first assertion follows from the definition of the cores.
For the second, sum \eqref{expansion} over $k\ge0$ using
$\sum_{k\ge r}\binom{k}{r}z^{k-r}=(1-z)^{-r-1}$,
and remove the $k=0$ term.
\end{proof}

\paragraph{Exact corroboration and scope.}
The accompanying standard-library verifier enumerates permutations
through size $9$, checks compression and reconstruction for
$113{,}979$ objects, and checks $2{,}808$ arbitrary primitive
insertions by direct pattern inspection. Two pattern detectors
agree on all $46{,}233$ permutations of sizes $1$ through $8$.
It also checks the coefficient identities and their binomial
simplification with exact integers. These computations corroborate
the proof; the all-parameter conclusion rests on
Lemma \ref{insertion} and the bijection above. The initial core
series are not evaluated in closed form here.

\section*{Task 2. Explicit initial series (proposed)}
Find closed forms or finite functional equations for $H_{3,1}$
and $H_{3,2}$. The corollary would then give an explicit expression
for $g_3(x,t)$, addressing the further question stated in the source.

\section*{Task 3. Sharpness of the component bound (proposed)}
Determine for which $a$ the series $H_{a,a-1}$ is nonzero and
find its least nonzero coefficient. This would decide when the
degree $a-1$ in the binomial expansion is necessary.

\section*{Task 4. Other forbidden patterns (proposed)}
Classify permutation classes whose primitive suffix blocks have
the same deletion and arbitrary-reinsertion property. Test this
first for the extra restrictions $1243$ and $2134$ studied in
\cite{GLW26}, and formulate the corresponding finite expansions.

\section*{Task 5. Asymptotics at fixed $a$ (proposed)}
Determine the dominant singularities of the finitely many
$H_{a,r}$, then use the displayed expression for $g_a$ to obtain
coefficient asymptotics and the limiting distribution of the
distance from $a$ to the maximum. No asymptotic claim is made here.

\paragraph{Status.}
Task 1 proves the full stated expansion conjecture. Tasks 2--5
are research proposals and carry no additional solution credit.
Review and computation were performed by the authoring agent;
no independent peer review or proof-assistant certification is claimed.

\begin{thebibliography}{9}\raggedright
\bibitem{GLW} J. B. Gil, O. A. Lopez, and M. D. Weiner,
\emph{A positional statistic for 1324-avoiding permutations},
Discrete Mathematics \& Theoretical Computer Science 26:1 (2024).
\href{https://doi.org/10.46298/dmtcs.12629}{doi:10.46298/dmtcs.12629}.
Current and journal-final text:
\href{https://arxiv.org/abs/2311.18227v3}{arXiv:2311.18227v3}.
\bibitem{GLW26} J. B. Gil, O. A. Lopez, and M. D. Weiner,
\emph{On separable permutations and three other pairs in the
Schr\"oder class},
\href{https://arxiv.org/abs/2603.25528v2}{arXiv:2603.25528v2} (2026).
\end{thebibliography}
\end{document}
